\section{Abóbora Contest}
Como meio didático e objeto de estudos sobre as maratonas de programação, serão utilizadas as questões da maratona do Abóbora Contest (Maratona própria do Instituto Federal Goiano Campus Rio Verde). Nesta seção serão apresentadas as questões presentes nas provas dos anos de 2022,2023,2024 e 2025. Incluindo, também, suas respectivas soluções. 

\begin{center}
 {\large {\textbf 
{\large \textbf{Informações Gerais}}
}}
\end{center}

Cada caderno contém $10$ problemas exclusivos criados por professores e alunos do Instituto Federal Goiano - Campus Rio Verde.

Segue abaixo as regras da respectiva maratona:

\begin{enumerate}
    \item \textbf{Sobre os nomes dos programas}
    \begin{enumerate}
        \item Para soluções em C/C++ e Python 3.8, o nome do arquivo-fonte não é significativo, pode ser qualquer nome desde que tenha as extensões .c, .cc, .py3.
        \item Para linguagem Java sua solução deve ser chamado de codigo\_do\_problema.java, onde codigo\_do\_problema é a letra maiúscula que identifica o problema. Lembre que em Java o nome da classe principal deve ser igual ao nome do arquivo.
    \end{enumerate}
    \item \textbf{Sobre a entrada}
    \begin{enumerate}
        \item A entrada de seu programa deve ser lida da entrada padrão.
        \item A entrada é composta de um único caso de teste, descrito em um número de linhas que depende do problema.
        \item Quando uma linha da entrada contém vários valores, estes são separados por um único espaço em branco; a entrada não contém nenhum outro espaço em branco.
        \item Cada linha, incluindo a última, contém exatamente um caractere final-de-linha.
        \item O final da entrada coincide com o final do arquivo.
    \end{enumerate}
    \item \textbf{Sobre a saída}
    \begin{enumerate}
        \item A saída de seu programa deve ser escrita na saída padrão.
        \item Quando uma linha da saída contém vários valores, estes devem ser separados por um único espaço em branco; a saída não deve conter nenhum outro espaço em branco.
        \item Cada linha, incluindo a última, deve conter exatamente um caractere final-de-linha.
    \end{enumerate}
\end{enumerate}

\newpage

%\include{Problema/abóbora22}

\newpage

%\include{ProblemaB_H/ProblemaB}

\newpage

%\include{ProblemaC_MB/ProblemaC}

\newpage

%\include{ProblemaD_AJ/ProblemaD}

\newpage

%\include{ProblemaE_DC/ProblemaE}

\newpage

%\include{ProblemaF_A/ProblemaF}

\newpage

%\include{ProblemaG_MB/ProblemaG}

\newpage

%\include{ProblemaH_H/ProblemaH}

\newpage

%\include{ProblemaI_AJ/ProblemaI}

\newpage

%\include{ProblemaJ_U/ProblemaJ}

\newpage
